\documentclass[12pt]{article}
\usepackage[utf8]{inputenc}
\usepackage[T1]{fontenc}
\usepackage[margin=1in]{geometry}
\usepackage{setspace}
\usepackage{xcolor}
\usepackage[colorlinks=true, urlcolor=blue!60!black, linkcolor=blue!60!black]{hyperref}
\setlength{\parindent}{2em}
\setlength{\parskip}{1.25em}
\renewcommand{\baselinestretch}{1.0}
\begin{document}
\begin{center}
{\LARGE \textit{Moore v. United States}, 602 U.S. 572 (2024)\par}
\vspace{0.8em}
{\large Case Brief --- Four-Page Edition\par}
\end{center}
\vspace{1em}

\section{Memory Jogger}

Taxpayers who owned 13\% of an Indian company challenged the 2017 Mandatory Repatriation Tax as an unapportioned tax on unrealized income; the Court upheld the tax 7--2. The majority held that Congress may attribute a pass-through entity's realized but undistributed income to its shareholders or partners without apportionment when the entity itself has not been taxed on that income, while reserving whether the Constitution requires realization (\href{https://supreme.justia.com/cases/federal/us/602/22-800/}{Justia}).

\section{Detailed Case Facts}

In 2006 Charles and Kathleen Moore invested \$40,000 for a 13\% stake in KisanKraft, an American-controlled foreign corporation (CFC) in India that supplies equipment to small farmers. KisanKraft was profitable every year through 2017 but reinvested its earnings, so neither it nor the Moores paid U.S. tax on them. Subpart F had long taxed U.S. shareholders of CFCs mainly on passive income. The Tax Cuts and Jobs Act of 2017 enacted the Mandatory Repatriation Tax (MRT), 26 U.S.C. \S\ 965, a one-time tax (8\% to 15.5\%) on U.S. shareholders owning at least 10\% of a CFC, attributing to them their pro rata shares of accumulated post-1986 earnings. The Moores' share was about \$508,000; they reported \$132,512 and paid \$14,729.

\section{Procedural History}

The Moores sued for a refund in the Western District of Washington, arguing that the MRT was an unapportioned direct tax and violated due process. The district court dismissed. The Ninth Circuit affirmed, declaring that ``realization of income is not a constitutional requirement,'' and denied rehearing en banc over Judge Bumatay's dissent. The Supreme Court granted certiorari on the Direct Tax Clause and Sixteenth Amendment question, heard argument December 5, 2023, and affirmed June 20, 2024.

\section{Judicial Votes}

Affirmed, \textbf{7--2}.
\begin{itemize}
\item \textbf{Majority (5):} Justice Kavanaugh, joined by Chief Justice Roberts and Justices Sotomayor, Kagan, and Jackson.
\item \textbf{Concurrences:} Justice Jackson; Justice Barrett, joined by Justice Alito, in the judgment.
\item \textbf{Dissent:} Justice Thomas, joined by Justice Gorsuch.
\end{itemize}

\section{Holding}

The MRT falls within Congress's taxing power. When an entity such as a partnership, S corporation, or CFC realizes income that stays undistributed and untaxed at the entity level, Congress may attribute that realized income to the owners and tax their pro rata shares without apportionment. The holding is limited to pass-through taxation; the Court reserved taxing both entity and owners on the same income, unapportioned taxes on wealth or net worth, and taxes on unrealized appreciation.

\section{Analysis of Opinions}

\textbf{Majority (Kavanaugh).} Congress's taxing power is ``exhaustive,'' reaching direct taxes, which must be apportioned, and indirect taxes, which must be uniform. Income taxes were indirect until \textit{Pollock} (1895); the Sixteenth Amendment restored taxes on ``incomes, from whatever source derived,'' to the unapportioned class (\textit{Brushaber}). Without deciding whether realization is required, the Court held that the MRT taxes realized income---KisanKraft's---attributed to its shareholders. Precedents from \textit{Burk-Waggoner} to \textit{Heiner v. Mellon} and \textit{National Grocery} let Congress tax an entity or its owners on undistributed income; \textit{Eisner v. Macomber}'s ``capital, not income'' remark was dicta. Congress has taxed owners on entity profits since 1864, and the Moores' ``ad hoc distinctions'' would threaten partnerships, S corporations, Subpart F, and more---a ``fiscal calamity.'' Due process bars arbitrary attribution.

\textbf{Jackson (concurring).} The Sixteenth Amendment contains no realization requirement, \textit{Macomber} has been confined to its facts, and the MRT could also be sustained as an excise; relief from unpopular taxes lies at the ballot box.

\textbf{Barrett, with Alito (in the judgment).} Sixteenth Amendment ``income'' requires realization, and the Moores realized nothing; the majority misreads its precedents to allow disregard of corporate form. She concurred only because the Moores conceded Subpart F's validity and failed to distinguish the MRT from it.

\textbf{Thomas, with Gorsuch (dissenting).} The Direct Tax Clause protects state tax bases, and the Amendment overruled only \textit{Pollock}'s conflation of income and source; distinguishing them requires realization. The MRT taxed mere stock ownership, and the majority's reservation of wealth taxes was ``cheap dicta.''

\section{Examples: Future Applications}

\begin{enumerate}
\item \textbf{Same-Side: Annual CFC Shareholder Tax.} Congress taxes 10\% U.S. shareholders each year on their shares of a CFC's realized, undistributed profits. Under \textit{Moore}, the attribution is valid without apportionment.
\item \textbf{Same-Side: Partners Taxed on Retained Profits.} A real-estate partnership keeps \$10 million in profits, and a minority partner who cannot force a distribution is taxed on his share. This is the paradigm case; \textit{Heiner v. Mellon} upheld it.
\item \textbf{Opposite-Side: Billionaire Appreciation Tax.} An unapportioned 2\% annual tax on unrealized appreciation of assets held by individuals worth over \$100 million. No entity realized income, so \textit{Moore} cannot save it; four Justices view realization as required.
\item \textbf{Opposite-Side: Taxing Customers on a Supplier's Profits.} Congress attributes a foreign manufacturer's profits to U.S. buyers of its products. \textit{Moore} allows attribution only to shareholders or partners, and due process bars arbitrary attribution (\textit{Burnet v. Wells}).
\item \textbf{Fence-Sitter: Pass-Through Tax on Public Shareholders.} Congress repeals the corporate tax for public companies and attributes profits to all shareholders. \textit{Moore}'s logic applies, but attribution to passive public shareholders may be arbitrary; footnote 8 left domestic and widely held corporations open.
\end{enumerate}

\section{Critique}

\textbf{Doctrinal minimalism versus consequentialism.} Critics argue that Justice Kavanaugh ducked the question on which certiorari was granted to avoid a ``fiscal calamity,'' leaving lower courts an undertheorized ``non-arbitrary attribution'' standard drawn from due process rather than a clear Sixteenth Amendment rule.

\textbf{Formalism versus economic substance.} The split between Barrett and Thomas, who insist on legal separation between corporation and shareholder, and Kavanaugh and Jackson, who treat entity-level realization as enough, reflects a recurring question whether Article I's tax limits track formal entities or economic control and revenue realities.

\section{Key Quotations}

\begin{quote}
``So the precise and narrow question that the Court addresses today is whether Congress may attribute an entity's realized and undistributed income to the entity's shareholders or partners, and then tax the shareholders or partners on their portions of that income. This Court's longstanding precedents, reflected in and reinforced by Congress's longstanding practice, establish that the answer is yes.''
\end{quote}
\noindent --- Justice Kavanaugh, Opinion of the Court.

\begin{quote}
``The Sixteenth Amendment's reference to income `derived' from any source encompasses a requirement that income, to be taxed without apportionment, must be realized.''
\end{quote}
\noindent --- Justice Barrett, concurring in the judgment.

\begin{quote}
``Sixteenth Amendment `incomes' include only income realized by the taxpayer. The text and history of the Amendment make clear that it requires a distinction between `income' and the `source' from which that income is `derived.' And, the only way to draw such a distinction is with a realization requirement.''
\end{quote}
\noindent --- Justice Thomas, dissenting.

\end{document}
